This study has several limitations. First, while we conducted experiments with twelve open-source LLMs, our findings may not be generalizable to other open-source models or newer architectures due to variations in design and training data. This limitation suggests that further research could explore a wider array of models to validate the applicability of our findings.

Second, our current repetition detection method focuses solely on non-word character and text repetition. This may overlook other relevant forms of repetition, such as word salad repetition, which could further impact the quality of generated content. Addressing these additional types of repetition in future work could lead to a more comprehensive understanding of repetition issues in LLMs.

Third, this research examines two specific tasks: question answering and machine translation. While these tasks are significant, extending the study to include other applications, such as text summarization, could provide deeper insights into LLM behavior concerning text repetition across various NLP tasks. This expansion would help in understanding how different tasks may influence repetition patterns and the overall performance of LLMs.

Finally, this work concentrates on the repetition penalty parameter (\rpp). However, other methods, such as sampling techniques, temperature adjustments, and context size modifications, can also be used to mitigate repetition issues in generative models. While \rap can be easily applied alongside these methods, further research evaluating \rap across different repetition control methods would enhance our understanding of how to manage repetition issues in LLMs more effectively.

% Our research has a few limitations worth noting. First, although we have conducted experiments in twelve open-source LLMs, it might still hard to have a generlizable conclusion for other open-source LLMs or newer models. Second, regarding repetition detection, currently we only consider non-word character and text repetition, however, it could have more type of repetition, such as, word salad repetition. Third, in this paper, we study two tasks of question answering and machine translation. Extending this study to other applications like text summarization could provide more understanding LLMs' behaviors in generating text repetition across different NLP tasks.

    % First, 
    % we focused on a limited selection of open-source LLMs, which might make it harder to generalize our findings to proprietary or newer models; expanding this scope in future studies could provide more insights. Second, while the repetition ratio (RR) is a cool new metric, we haven’t fully explored how it stacks up against more established ones, so future work should look into how RR correlates with and maybe even enhances other metrics in different linguistic contexts. Lastly, adding our repetition detection algorithm could introduce some computational overhead, which might affect the performance and scalability of LLMs. Future versions of the framework will need to find a good balance between complexity and efficiency for wider use.
    % \item \textbf{Scope of LLMs Tested}: Our study was limited to a selection of open-source LLMs, potentially constraining the generalizability of our findings to proprietary or newly developed language models. Future research could expand this scope to assess the framework’s effectiveness across a broader range of LLMs.
    % \item \textbf{Metric Dependence}: While the repetition ratio (RR) is an innovative addition to performance metrics for LLMs, its effectiveness and reliability relative to more established metrics remain underexplored. Future studies should examine how RR correlates with and possibly enhances other metrics across various linguistic contexts.
    % \item \textbf{Complexity and Performance}: The integration of our repetition detection algorithm may introduce computational overhead, potentially impacting the performance and scalability of LLMs. Future iterations of the framework will need to balance complexity with efficiency to ensure broader applicability.
    % \item \textbf{Adjustment of Repetition Penalties}: Achieving an optimal balance of the repetition penalty parameter (RPP) without manual intervention remains a challenge. This process may benefit from more sophisticated, AI-driven automation techniques in future developments.
    % \item \textbf{Task Specificity}: Our research focused primarily on question-answering tasks. Extending the framework to other applications like summarization could provide a more comprehensive understanding of its utility across different NLP tasks.
    % \item \textbf{False Positives in Repetition Detection}: The current repetition detection algorithm may produce minor false positives when identifying closely related words, such as in the phrase “resolve solve.” Adjusting the threshold for text repetition detection could mitigate this issue. We plan to address this refinement in our future work to enhance accuracy.